\paragraph{Inspection-based UX evaluation.}
Cognitive walkthroughs are a classic, theory-driven inspection method for identifying usability issues without running user studies~\cite{polson1992cognitive}.
At scale, organizations adopt metrics frameworks such as HEART to connect product goals to measurable signals and metrics~\cite{rodden2010heart}.
The System Usability Scale (SUS) remains a widely used, lightweight instrument for perceived usability~\cite{brooke1996sus}.

\paragraph{Mining user feedback at scale.}
App reviews and related feedback sources have been extensively studied in requirements engineering and software maintenance,
with systematic surveys cataloging use cases and evaluation gaps~\cite{dabrowski2022review}.
Public corpora of reviews, support conversations, and issues enable reproducible studies, but they rarely include
task completion times or controlled experimental conditions.
Recent datasets focus on structured taxonomies of end-user issues, including low-rated Amazon Appstore reviews with a labeled subset~\cite{amazonlowrated2026dataset}.

\paragraph{LLMs as simulators and the ``human surrogate'' debate.}
Recent work explores ``silicon samples'' and LLM-based simulation of human responses~\cite{argyle2023outofone,aher2023simulate,horton2023homosilicus}.
In HCI, LLM-powered agents and personas have been proposed for scalable simulation and ideation~\cite{park2023generativeagents,choi2025proxona,schuller2024personas,prpa2024syntheticpersonae}.
However, multiple studies highlight validity risks and misalignment between LLM outputs and human distributions~\cite{bisbee2024perils,amirova2024algorithmicfidelity,gao2024scylla,li2025personacatch}.
We position our work as \emph{decision support for early iteration}, not a replacement for human studies:
we use proxy validation to detect egregious mismatches and to guide prompt/taxonomy iteration,
while reserving stronger paid-model calibration for the final evaluation pass.

\paragraph{LLM-based evaluation and rubric reliability.}
LLM-as-a-judge methods are increasingly used as scalable approximations to human evaluation~\cite{liu2023geval,zheng2023mtbench},
but they can exhibit systematic biases (e.g., position and verbosity effects) and inconsistent rubric application~\cite{wang2024notfair,shi2024positionbias,stureborg2024inconsistent}.
We therefore emphasize (i) schema-constrained outputs, (ii) decision trails tied to concrete artifacts, and (iii) small gold-labeled calibration sets when available.

\paragraph{Positioning.}
Unlike prior LLM simulation work that targets individual-level behavioral fidelity or post-hoc user-study surrogates, our contribution is a \emph{protocol} layer: a lightweight proxy-validation loop that quantifies distributional alignment between simulation outputs and public corpora, with reproducible alignment metrics, bootstrap confidence intervals, and artifact-first reporting.
This complements existing frameworks by providing structured go/no-go signals for prompt and taxonomy iteration \emph{before} committing to expensive human studies or large-scale model calibration.
